\section{Local Markers}

In Chapter \ref{chap:local_markers}, we considered the role of band topology in insulators where the presence of disorder makes the notion of band structure ill-defined. In particular, we asked whether it was possible to construct a topological invariant analogous to the Chern number, that could be robustly formulated without a notion of k-space. We provided a review of three approaches to this problem, each proposing a different way to build a local `marker' of Chern number \cite{bianco_mapping_2011,loring_guide_2019,kitaev_anyons_2006}, and discussed the ways that these quantities are related to one another. In each of these proposals, the local Chern number was constructed without a clear understanding of whether these may be measured as a real, physical observable. This approach has two drawbacks. Firstly, the markers are of little use to an experimentalist, since there is no obvious physical observable to which they correspond. Secondly, although the markers are useful theoretical tools, they do not provide a satisfying intuitive answer to the question of what it meant for a material to have a physical notion of local topology. In \textsection\ref{sec:physical_interpretation_kitaev_marker} we proposed a solution to this problem by constructing a physical observable -- a localised Hall conductivity -- that corresponds exactly to Kitaev's local Chern marker \cite{kitaev_anyons_2006} when evaluated in the adiabatic limit. We provided a careful discussion of the exact operators that must be measured to give a quantised topological result, in particular showing that an extremely weak notion of local conductivity is required, where we must be able to differentiate between conductivity that is localised around a point deep within the bulk of the system and the conductivity that is transported along the edge modes, with no current in the in-between area. Nevertheless, many open questions remain.

Firstly, the arguments used to derive this local observable are essentially a prescription for measuring the quantity in the lab: a step-function potential is used to excite current across the system in the vertical direction. The cross-conductance is then determined by measuring the horizontal transport of electrons across a line dividing the system. Thus, an obvious next step will be to look at experimentally verifying the proposed method. In particular, note that the observable we present benefits from many of the remarkably invariant properties of Kitaev's marker, which we discussed in \textsection\ref{sec:kitaev_marker}. The step function potential can be smoothed into a localised slope, and the two lines -- one for an electric field and one for a measurement -- need not be exactly perpendicular.

From a theory perspective, this proposal also suggests a number of avenues for further study. One particularly interesting example comes when considering how these markers may be extended to strongly interacting matter, especially since the Chern number was shown by Niu, Thouless and Wu to have a straightforward many-body construction as an integral over twisted boundary conditions \cite{niu_quantized_1985}. In general, local markers have been constructed in a purely non-interacting context, as a property of the projector onto a band of single-particle states. However, in any interacting system, the ground state corresponds to a single highly complex state, and it is usually impossible to separate this into single-particle states. Thus, one cannot use the existing constructions of these local markers, although some progress has been made using Green's functions \cite{markov_local_2021}. Although the marker itself does not work in a many-body context, the intuitive picture we have presented -- that of currents and fields -- is still applicable. An interesting avenue of further study is to test whether it is possible to exploit this physical picture in order to construct a local Chern marker that is robust in a many-body framework. In particular, the Kitaev honeycomb model (and Kitaev materials in general) presents an interesting challenge, since the signature of the non-Abelian phase in the Kitaev model is a quantised \emph{thermal} Hall effect. It remains an open question whether a similar reasoning can be used to define a localised measure of thermal conductivity \cite{luttinger_theory_1964} that was robust even away from the exactly solvable Kitaev point.

Finally, this thesis has focussed exclusively on two-dimensional phases with no special symmetries -- those characterised by the Chern number. However, this is a small window into a much larger classification, that of symmetry-protected topological (SPT) phases \cite{ryu_topological_2010}. Depending on dimensionality and symmetry, a broader range of topological phases can be constructed with qualitatively different invariants, for example winding numbers in the Su-Schrieffer-Heeger model \cite{su_solitons_1979}, $\mathbb Z_2$-type invariants found on the Kitaev Chain \cite{kitaev_unpaired_2001}, or the Bernevig-Hughes-Zhang model \cite{bernevig_quantum_2006}. In general, local markers have been constructed for many of these phases \cite{chen_universal_2023,huang_theory_2018,sykes_local_2021,hannukainen_local_2022}, however few have been constructed using Kitaev's original formalism, and even fewer with a robust physical interpretation. Thus, a valuable extension would be to look at extending the arguments presented here to across the full classification of SPT phases.

